% Results (Task 46, Scientific Reports): v3 Results with statement subheadings; each subsection opens with the
% one-sentence answer to its research question. Every result number is a macro. All caveats of plan 42 section 3.2
% are in this section.
\section*{Results}
\label{sec:results}

Unless stated otherwise, numbers are values on the primary row set V1, from which test rows with a near duplicate in
the training set were removed (Methods), on the \nCommonClasses{} common classes at the budget
$\epsilon = \paramEpsHead$, given as the mean over \nSeeds{} seeds with the seed range in brackets.

\subsection*{Attention changed little under constrained attacks but was less faithful than integrated gradients and
occlusion}
\label{sec:rq1}
Rollout attention agreed weakly to moderately with attribution methods and was less faithful than integrated gradients
(IG) and occlusion on both datasets; it changed little under constrained attacks that targeted it, and after an evasion
its faithfulness fell together with that of the other explanation methods.

On clean data, the Spearman correlation of rollout attention with the attribution methods was highest with
GradientSHAP (\spRolloutGShapVoneSeventeen{} on 2017, \spRolloutGShapVoneEighteen{} on 2018), about one third with IG
(\spRolloutIGVoneSeventeen{}, range \spRolloutIGVoneSeventeenRange; \spRolloutIGVoneEighteen{}, range
\spRolloutIGVoneEighteenRange), and lowest with TreeSHAP of the XGBoost model (\spRolloutTreeShapVoneSeventeen{} and
\spRolloutTreeShapVoneEighteen; Table~\mainTabClean).
Last-layer attention showed almost no agreement (Spearman with IG \spLastIGVoneSeventeen{} and \spLastIGVoneEighteen).
Rollout was more faithful than a random ranking and less faithful than IG and occlusion: its deletion area under the
curve (AUC, lower is more faithful) was \aucDelRolloutVoneSeventeen{} on 2017 and \aucDelRolloutVoneEighteen{} on 2018,
against \aucDelIGVoneSeventeen{} and \aucDelIGVoneEighteen{} for IG, \aucDelOcclVoneSeventeen{} and
\aucDelOcclVoneEighteen{} for occlusion, and \aucDelRandomVoneSeventeen{} and \aucDelRandomVoneEighteen{} for random,
and the insertion AUC gave the same picture.
Among the \nMethods{} methods, rollout ranked \rankDelRolloutSeventeen{} and \rankDelRolloutEighteen{} by deletion AUC,
and the ordering of the methods was similar across the datasets (Spearman correlation between the two orderings
\ordSpDelVone{} for deletion and \ordSpInsVone{} for insertion AUC).

Under the constrained attack A2, which kept the prediction and minimised the rollout attention on the clean
top-\paramTopK{} IG features, the attention mass on these features kept \massAtwoVoneSeventeen{}
(\massAtwoVoneSeventeenRange) of its clean value on 2017 and \massAtwoVoneEighteen{} (\massAtwoVoneEighteenRange) on
2018 (Table~\mainTabAttn, panel a; Supplementary Fig.~\suppFigMassRatio).
Under the random control A0 at the same budget it kept \massAzeroVoneSeventeen{} and \massAzeroVoneEighteen, and the
joint evasion and attention attack A3 reached \massAthreeVoneSeventeen{} and \massAthreeVoneEighteen.
An unconstrained reference moved the mass further (\massUncVoneSeventeen{} and \massUncVoneEighteen), with rows that
are not valid flows and a budget that is not in the units of the constrained attacks.
Attention was therefore stiff under the feature constraints; this does not show that attention cannot be manipulated,
and IG was not attacked at an equal budget.
For context, the evasion attack A1 succeeded on \succAoneVoneSeventeen{} and \succAoneVoneEighteen{} of the source
rows, \aoneMinusAzeroVoneSeventeen{} and \aoneMinusAzeroVoneEighteen{} more than A0.

On rows whose prediction A1 flipped to benign, the margin AOPC gap (the gap between most- and least-relevant-first
removal) fell after subtracting the change under A0 for IG (\flipIGVoneSeventeen{} and \flipIGVoneEighteen),
occlusion (\flipOcclVoneSeventeen{} and \flipOcclVoneEighteen), rollout (\flipRolloutVoneSeventeen{} and
\flipRolloutVoneEighteen), the static class-mean rollout ranking (\flipStaticRolloutVoneSeventeen{} and
\flipStaticRolloutVoneEighteen) and $|x|$, while the random ranking stayed near zero (Table~\mainTabAttn, panel b).
The last-layer attention changed with a sign that varied across seeds: \flipLastVoneSeventeen{} on 2017 (seed range
\flipLastVoneSeventeenRange) and \flipLastVoneEighteen{} on 2018 (\flipLastVoneEighteenRange).
On A1 rows that did not flip the change was small (IG \notflipIGVoneSeventeen{} and \notflipIGVoneEighteen), and on
rows that A0 flipped by chance the loss was of the same size as on flipped A1 rows on 2017 and larger on 2018 (IG
\azeroflipIGVoneSeventeen{} and \azeroflipIGVoneEighteen).
Because the loss appeared also for static rankings that the attack cannot change and on rows flipped by a random
perturbation, it is consistent with a property of the model at the adversarial point rather than with manipulation of
attention; we did not test the mechanism.
Three smaller 2017 results did not hold in the same direction on 2018: the advantage of a static class-mean rollout
ranking over the per-row ranking, the advantage of $|x|$ over rollout, and the advantage of a single attention head
selected on validation data (Supplementary Note~\suppNoteResults); a parameter-randomisation check showed that the
attention explanations changed when the learned weights were replaced (Supplementary Note~\suppNoteRand).

\subsection*{Consistency scores raised evasion against exact-objective attackers}
\label{sec:rq2}
The consistency scores helped against an attacker that ignored them, largely without attention, but against attackers
that optimised the detector's exact score the extended detector left more attack rows successful and undetected than
the base detector on both datasets, and the gain measured with a stand-in attacker did not hold.

The base detector D combines confidence and distance scores; D$_+$ adds \nConsistency{} attention-attribution
consistency scores (Methods).
Both detectors separated their validation positives well (validation AUROC of D \valAurocDSeventeenRange{} on 2017
and \valAurocDEighteenRange{} on 2018 over seeds, D$_+$ \valAurocDplusSeventeenRange{} and
\valAurocDplusEighteenRange); these values are in-sample, because the flipped validation positives were also used to fit
the detectors, while the realised test FPR of D, \testFprDSeventeenRange{} and \testFprDEighteenRange, is the
out-of-sample figure.
Against A6, which ignores the consistency scores, D$_+$ detected more of the evading rows than D: the paired difference
in true-positive rate (TPR) D$_+$ minus D on A6 rows ($\lambda = 1$) was \tprAsixOneVoneSeventeen{}
(\tprAsixOneVoneSeventeenCI) on 2017 and \tprAsixOneVoneEighteen{} (\tprAsixOneVoneEighteenCI) on 2018
(Table~\mainTabAttacks).
The control detectors show that this gain was largely obtained without attention (Table~\mainTabControls): the TPR
difference to D was \ablBMinusDSeventeen{} and \ablBMinusDEighteen{} for the attribution-only control b and
\ablCtwoMinusDSeventeen{} and \ablCtwoMinusDEighteen{} for c2 with logit and distance statistics, while c0
(\ablCzeroMinusDSeventeen) and c1 (\ablConeMinusDSeventeen) show that the regulariser and the number of scores do not
explain it.
Shuffled attention from another row of the same class (a1) kept \ablAoneMinusDSeventeen{} and
\ablAoneMinusDEighteen, attention from clean benign rows (a2) kept \ablAtwoMinusDSeventeen{} and
\ablAtwoMinusDEighteen, and the pairing of a row's own attention with its attribution added \ablPairSeventeen{}
(\ablPairSeventeenCI) on 2017 and nothing detectable on 2018 (\ablPairEighteen, \ablPairEighteenCI).

Against A8, which models the consistency scores through differentiable stand-ins, the gain largely disappeared (TPR
difference \tprAeightOneVoneSeventeen{} on 2017 and \tprAeightOneVoneEighteen{} on 2018).
Against A8x, which optimises the exact score of D$_+$ wherever an exact gradient exists, D$_+$ detected fewer of the
attack rows than D (\xtprAxOneVoneSeventeen{}, \xtprAxOneVoneSeventeenCI, on 2017 and \xtprAxOneVoneEighteen{} on
2018; Table~\mainTabAttacks), and the difference was negative for every setting and row set (largest value
\xtprAxMaxSeventeen{} on 2017; Supplementary Table~\suppTabAxTpr).
A8x flipped about as many rows as A8 (success \xsuccAxOneVoneSeventeen{} against \xsuccAeightOneVoneSeventeen{} on
2017) but left more of them undetected by D$_+$ (\xundetAxOneVoneSeventeen{} against \xundetAeightOneVoneSeventeen).
These rows were optimised against D$_+$, so this comparison alone is not evidence; the equal-effort union is.

Under equal attacker effort, with D facing A1 and the three settings of A6 and D$_+$ facing A1 and the three settings
of A8x, D$_+$ left more source rows successful and undetected than D (Fig.~\mainFigHeadline, Table~\mainTabUnion):
\axHeadVoneSeventeenDplus{} against \axHeadVoneSeventeenD{} on 2017, a difference of \axHeadVoneSeventeen{}
(\paramCILevel{} confidence interval (CI) \axHeadVoneSeventeenCI), and \axHeadVoneEighteenDplus{} against
\axHeadVoneEighteenD{} on 2018, a difference of \axHeadVoneEighteen{} (\axHeadVoneEighteenCI).
With A8 instead of A8x the 2017 difference had been \axStandinVoneSeventeen{} (\axStandinVoneSeventeenCI), and on 2018
\axStandinVoneEighteen{} (\axStandinVoneEighteenCI).
Per seed, the differences were \axHeadVoneSeventeenSeeds{} on 2017 and \axHeadVoneEighteenSeeds{} on 2018
(seeds 0, 1 and 2), with \axHeadOppSeedsPhrase{} of the opposite sign.
By the outcome rule recorded before the A8x runs (Methods), the 2017 gain of D$_+$ against A8 did not hold against the
exact-objective attacker.
On 2017 the direction held for each attack setting (\axLamOneVoneSeventeen, \axLamHundredVoneSeventeen{} and
\axNormVoneSeventeen; Supplementary Table~\suppTabAxSetting) and on the other row sets (Table~\mainTabNeardup).
With the A8 and A8x candidates together against D$_+$, an unequal-effort union that favours the attacker of D$_+$, the
2017 difference was \axSensVoneSeventeen{} (\axSensVoneSeventeenCI); it is a sensitivity analysis, not the headline.

These numbers come with caveats.
They rest on \nSeeds{} seeds, with \axHeadOppSeedsPhrase{} of the opposite sign, and the intervals are bootstraps over
source rows within seed that do not capture variation between seeds.
A8x still uses surrogate gradients for the \nRankScores{} rank-based scores; a better search on these terms could only
find further undetected rows for D$_+$.
On 2018 D alone left only \axHeadVoneEighteenD{} of the source rows successful and undetected, so the 2018 difference
sits near a floor and its lower bound is close to $0$.
The equal-effort protocol fixes candidates, steps and settings, not optimiser quality, and whether A6 is as strong an
attack on D as A8x is on D$_+$ was not measured.

A8x keeps the iterate with the lowest true D$_+$ score, while A6 keeps its own rule, so part of the reversal could have
come from the selection rule rather than from the consistency scores.
The attacker A6x gives D's attacker the selection rule of A8x.
It kept a different iterate in \sxRowsDiffer{} of \sxRowsAll{} attacked rows but changed the flipped status of
\sxFlipChanged{} rows and the undetected status of \sxUndetChanged{} rows, because A6 already ranks iterates that are
flipped and below the threshold under the exact score of D first.
D's undetected share under A1 and A6x therefore equalled that under A1 and A6 in every seed and row set (largest
absolute difference \sxSelMaxAbs).
With A6x the 2017 difference D$_+$ minus D kept a lower bound of \sxHeadVoneSeventeenLo{} (\sxHeadVoneSeventeenCI;
Fig.~\mainFigHeadline), and by the rule recorded before the A6x runs the reversal does not come from the selection
rule.
A6x matches the selection rule only, not the optimiser: the A8x rows were not aimed at D, yet D left a similar share of
them undetected as of the A6 rows aimed at it (2017: \evadeDAxOneSeventeen{} against \evadeDAsixOneSeventeen{} for
$\lambda = 1$), so the relative optimiser strength remains unmeasured.

\nControlsAttacked{} of the \nControlsAll{} controls were attacked by adversaries that modelled their own scores
(Table~\mainTabUnion).
Against D, the controls with exact attackers stayed close to D on 2017 (c0 \cmdCzeroVoneSeventeen, c1
\cmdConeVoneSeventeen, c2 \cmdCtwoVoneSeventeen), while a1 left more rows undetected (\cmdAoneVoneSeventeen,
\cmdAoneVoneSeventeenCI); on 2018, c0 (\cmdCzeroVoneEighteen) and a1 (\cmdAoneVoneEighteen) left fewer rows
undetected than D and c2 more (\cmdCtwoVoneEighteen), so no control left clearly fewer rows undetected than D on both
datasets.
Against D$_+$ under A8x, every control except a1 on 2017 and c2 on 2018 left fewer rows undetected (2017: c0
\cmxCzeroVoneSeventeen, c1 \cmxConeVoneSeventeen, c2 \cmxCtwoVoneSeventeen, a1 \cmxAoneVoneSeventeen{}
\cmxAoneVoneSeventeenCI, a2 \cmxAtwoVoneSeventeen; 2018: c2 \cmxCtwoVoneEighteen{} \cmxCtwoVoneEighteenCI).
For c0, c1 and c2 this comparison is at roughly equal attacker strength: exact against exact apart from the surrogate
ranks of A8x, and their attackers keep the selection rule of A6, which A6x showed to be equivalent for D but which was
not tested for them.
The controls a1 and a2 faced stand-in attackers (median $R^2$ of the stand-ins \rtwoMedAoneSeventeen{} and
\rtwoMedAtwoSeventeen{} on 2017), so their comparison with D$_+$ under A8x is not at equal attacker strength and their
shares may be too low, as A8's were for D$_+$.
Control b was not attacked adaptively, because its attack run exceeded the memory of the graphics card.
The comparisons with D$_+$ under the stand-in attacker A8 are in Supplementary Table~\suppTabUnionAeight.

\subsection*{The datasets differed more than the detectors}
\label{sec:rq3}
The level of successful and undetected rows depended strongly on the dataset and was dominated by one class on each,
the sign of D$_+$ minus D under A8x was the same on both datasets but its size was not, and removing near duplicates
changed no conclusion.

Under A6 with $\lambda = 1$, the share of source rows that evaded D undetected was \undetDAsixOneVoneSeventeen{} on
2017, with a seed range of \undetDAsixOneVoneSeventeenRange.
On 2018 it was \undetDAsixOneVoneEighteen, with a seed range of \undetDAsixOneVoneEighteenRange.
Under the equal-effort union the shares were \axHeadVoneSeventeenD{} and \axHeadVoneEighteenD.
On 2018, \pcTopClassEighteen{} was the only class with a non-trivial share (\pcTopEighteen), and every other class had
at most \pcRestMaxEighteen.
On 2017, \pcTopClassSeventeen{} dominated (\pcTopSeventeen), and every other class had at most \pcRestMaxSeventeen{}
(Supplementary Table~\suppTabPerClassVone).
Under A8x against D$_+$, DDoS again dominated on 2018, while on 2017 SSH-Patator joined Bot, with a seed range that
spans almost the whole scale (Supplementary Table~\suppTabPerClassVone).
The difference between the datasets in detector strength was larger than the difference between D and D$_+$.

Test rows with a near duplicate in the training set were common in both datasets (\ndAllSeventeen{} and
\ndAllEighteen{} of all test rows) and concentrated in DoS Slowloris and, on 2017, web attack.
On 2017, flagged rows were easier for both detectors: on A6 rows ($\lambda = 1$) the TPR of D was \tprDNdSeventeen{} on
flagged and \tprDOtherSeventeen{} on other source rows, and the FPR of D was \fprDNdSeventeen{} on flagged and
\fprDOtherSeventeen{} on other benign test negatives (Supplementary Note~\suppNoteNeardup).
Moving from all source rows (V0) to V1 and to V2, which also removes DoS Slowloris and web attack
(Table~\mainTabNeardup), changed the attention mass ratio under A2 and A3 by at most \sensMaxMassChange{} and the
faithfulness loss on flipped rows for IG and rollout by at most \sensMaxFlipChange.
The 2017 equal-effort difference under A8x was \axHeadVzeroSeventeen{}, \axHeadVoneSeventeen{} and
\axHeadVtwoSeventeen{} on V0, V1 and V2, and on 2018, where V1 equals V0, \axHeadVoneEighteen{} and
\axHeadVtwoEighteen{} on V2.
With the stand-in attacker A8 the 2017 difference changed more with the row set (\axStandinVzeroSeventeen,
\axStandinVoneSeventeen, \axStandinVtwoSeventeen).
